\subsection{The finite presentation of the global spectrum}
\label{sec:annex-p10-corrected-global-spectrum}

\begin{annexitem}{annex:p10-spectrum-lattice}{thm:p10-spectrum-lattice}{Finite lattice construction and change of base}
\begin{proof}
\emph{Finite joins without a natural numbers object.}
The definition of \(K(X)\) by intersection gives the following induction
principle: a property of its members holds for all of them if it holds
for the empty subobject and singletons and is closed under binary union.
Images, unions, and products of Kuratowski-finite subobjects are
Kuratowski-finite, by this induction.

If \(M\) is a join semilattice with bottom and \(f:X\to M\),
every \(U\in K(X)\) has a unique least upper bound of its \(f\)-image.
For the empty subobject use the bottom, for a singleton use its value,
and for a union use the binary join of the two bounds.
The least-upper-bound characterization gives uniqueness and independence
from any construction of \(U\).
The internal unique-existence statement defines an arrow.
It makes \(K(X)\), with union and empty subobject, the free
join semilattice with bottom on \(X\).
The dual argument supplies finite meets in a meet semilattice with top.

\emph{The free bounded distributive lattice.}
Use \(K(K(X))\) as a set of disjunctive normal forms.
A member \(U\) represents a finite join of finite meets:
its members are clauses, each a finite subobject of \(X\).
Define a preorder by
\[
 U\preceq V
 \quad\Longleftrightarrow\quad
 \forall s\in U\ \exists t\in V\ (t\subseteq s).
\]
Reflexivity uses \(t=s\); transitivity composes the two inclusions.
Join is union of clause families. Meet sends \(U,V\) to
\(\{s\cup t:s\in U,\ t\in V\}\).
The bottom is the empty family and the top is the singleton
family containing the empty clause.
These are again finite subobjects by the preceding induction.

The preorder is preserved by both operations. For meet, a contained
clause from each factor gives their contained union; for join,
use the factor from which a clause came.
Pass to the effective quotient by mutual comparison.
Union and union of clauses give associativity and commutativity.
The empty family and empty clause give the bound laws.
A clause \(s\cup t\) contains \(s\), giving one absorption
inequality; for the other, repeat \(s\) in both factors.
Distribution follows by distributing the two possible sources of
a clause in a union. Thus the quotient is a bounded distributive lattice.

Send \(x\) to the single clause \(\{x\}\).
For any map \(f:X\to M\) into a bounded distributive lattice,
send \(U\) to
\[
 \bigvee_{s\in U}\ \bigwedge_{x\in s}f(x).
\]
The finite bounds exist by the first paragraph.
If \(t\subseteq s\), the meet over \(s\) is below the meet over
\(t\), so this formula respects the preorder and the quotient.
It preserves join and meet, the latter by finite distributivity,
proved by induction on the two finite clause families.
It preserves the bounds and extends \(f\).
Every clause and clause family is its indicated finite meet and join
of generators, so no other extension is possible.

\emph{The ring relations.}
Apply the free construction to \(X=A\).
Intersect all lattice congruences containing the displayed ring
relations, writing each inequality as the equation
\(x\wedge y=x\).
The total congruence is one such congruence, so the intersection
is defined. It is an equivalence relation compatible with all
lattice operations. Its effective quotient is \(L_A\).
A map from \(A\) into a bounded distributive lattice extends to
the free lattice and factors through this quotient exactly when
the four required relations hold. This proves its universal property.

\emph{Preservation by inverse image.}
For any variety defined by finitary operations and equations,
a finite-product-preserving inverse-image functor \(f^*\)
transports its algebra objects. The right adjoint \(f_*\)
preserves finite products and transports the same algebra
structures in the reverse direction.
The ordinary adjunction restricts to algebra homomorphisms:
the adjoint transpose of a homomorphism preserves an operation
because the unit, counit, and finite-product comparisons commute
with that operation. The same verification applies to constants.
Thus the transported algebra functors are adjoint.

In particular \(f^*\) preserves free bounded distributive lattices,
as follows by applying that adjunction to their defining hom-set
bijections.
The same argument applied to join semilattices with bottom gives
\(f^*K(X)\cong K(f^*X)\). Thus the finite subobjects used to index
basic covers are transported to the corresponding finite subobjects,
rather than assumed to be preserved as arbitrary power objects. It also preserves their coequalizers.
The ring relations above can be presented as two parallel
homomorphisms from a free lattice on their object of instances
to the free lattice on \(A\); the quotient is their coequalizer.
The instance object uses only finite coproducts and finite
products of \(A\), and its two arrows use the ring operations.
They are transported to the corresponding instance object and
arrows for \(f^*A\).
Their coequalizers therefore give the claimed canonical isomorphism
\(f^*L_A\cong L_{f^*A}\).
This argument concerns finitely presented algebraic structure;
it does not assert that \(f^*\) preserves arbitrary power objects.

\emph{Its ideal frame.}
An ideal of a bounded distributive lattice is downward closed,
contains zero, and is closed under finite joins.
Intersections give meets; the least ideal containing a union gives
joins. For an ideal \(K\) and an ideal \(R\), the subobject
\[
 (K:R)=\{a:\ \forall r\in R,\ a\wedge r\in K\}
\]
is again an ideal, using distributivity for join closure.
If \(K=\bigvee_j(R\cap I_j)\), then every \(I_j\) lies in
\((K:R)\). Hence their join lies there.
An element \(a\) of \(R\cap\bigvee_j I_j\) consequently satisfies
\(a=a\wedge a\in K\). The reverse inclusion is immediate.
This proves the frame distributivity law.

Every ideal is the join of the principal ideals of its members.
For a bounded lattice homomorphism \(e:L_A\to H\) into a frame,
the unique extending frame map on ideals is
\(I\mapsto\bigvee_{u\in I}e(u)\).
To verify arbitrary joins, use the ideal
\(\{u:e(u)\leq v\}\) for a proposed upper bound \(v\).
To verify binary meets, distribute the two joins and use
\(e(u)\wedge e(v)=e(u\wedge v)\).
The top is preserved because \(e(1)=1\).
These verifications also prove uniqueness by principal generation.

By the presentation of \(L_A\), this ideal frame and the
radical-ideal frame of
\cref{thm:p10-radical-spectrum-presentation}
have the same universal property on generators \(A\).
The unique generator-preserving frame maps in both directions
are inverse: their composites fix every generator, and hence
every join of generators.
This gives the asserted canonical isomorphism.
\end{proof}
\end{annexitem}

\begin{annexitem}{annex:p10-spectrum-site-classification}{thm:p10-spectrum-site-classification}{Site}
\begin{proof}
Write \(L=L_A\).

\emph{The coverage and its sheaves.}
The singleton \(\{u\}\) covers \(u\).
Pulling a basic cover \(V\) of \(u\) back to \(w\leq u\)
gives \(\{w\wedge v:v\in V\}\).
It is finite, and its join is \(w\), by distributivity.
For transitivity, refine each member of a finite cover by a
finite cover. A finite union of the refinements has the required
join. This assertion requires no choice function:
induction on the finite covering subobject proves the existence
of a finite joint refinement, starting with the empty family
and combining two existentially supplied refinements by union.
Thus sieves containing basic covers satisfy the Grothendieck
topology axioms.

On the internal presheaf classifier, the associated local operator
sends a sieve to the sieve of positions where it contains a basic
cover. Inflationarity uses singleton covers; preservation of binary
intersections uses pairwise meets of two covers; idempotence is
the transitivity just proved. The internal-presheaf topos and its
left-exact sheaf reflector are supplied by the preceding sheaf
construction. Write that reflector as \(a\).
The constant-presheaf functor preserves finite limits and has
a right adjoint given by compatible dependent families of sections.
Composing it with \(a\) gives the inverse image \(p^*\) of
a geometric projection to the base.

Let \(y(u)\) be the representable presheaf for \(u\).
It is subterminal. It is a sheaf: if every member of a finite
cover of \(v\) lies below \(u\), their join \(v\) lies below \(u\).
The subterminals \(y(u)\) preserve finite meets and the top.
If \(V\) is a basic cover of \(u\), the family \(y(v)\to y(u)\)
for \(v\in V\) is jointly epimorphic in sheaves: its image
contains a covering family at every position in \(y(u)\).
In particular finite joins of these subterminals are the
representables of the corresponding joins in \(L\).
The representable \(y(0)\) is initial, since the empty cover
of zero forces its sole possible section in every sheaf.

\emph{A morphism from a lattice-valued predicate.}
Fix \(f:\mathcal F\to\mathcal E\) and a bounded lattice
homomorphism \(\chi:f^*L\to\Omega_{\mathcal F}\).
Let \(U_u\) be the subterminal classified by \(\chi(u)\).
They satisfy \(U_1=1\), \(U_{u\wedge v}=U_u\cap U_v\),
and \(U_{u\vee v}=U_u\vee U_v\).

First work with presheaves \(P\) on \(f^*L\) in \(\mathcal F\).
Form the object of representatives \((u,x)\), where
\(x\in P_u\) and \(U_u\) holds. Declare two representatives
\((u,x)\) and \((v,z)\) equivalent when locally there is
\(w\leq u,v\) with \(U_w\) and
\(x|w=z|w\).
Reflexivity uses \(w=u\); symmetry exchanges the two
representatives; transitivity uses the meet of two refinements.
These are internal existential statements. The equivalence
relation has an effective quotient, denoted \(G(P)\).

A map \(G(P)\to Y\) is exactly a compatible family of maps
\(P_u\times U_u\to Y\).
Currying identifies this with a presheaf map from \(P\) to
the presheaf with fibers \(Y^{U_u}\) and restriction by
precomposition. This gives an actual right adjoint to \(G\).
The representative quotient is also the coequalizer of the
two restriction maps, so no externally indexed colimit is
being assumed.

The functor \(G\) preserves finite limits.
For the terminal presheaf, \(u=1\) supplies a representative
everywhere, and any two representatives agree at their meet.
For binary products, representatives at \(u,v\) refine to
\(u\wedge v\); this supplies a representative of the pair.
Equality of pairs is likewise checked after a common refinement.
For an equalizer, equality of the two images of a germ means
that they agree on some supporting refinement. The restricted
representative then lies in the presheaf equalizer.
In each case these arguments establish an epimorphism and a
monomorphism for the canonical comparison, hence an isomorphism.
They use local representatives, not choices of global ones.

Pulling an \(\mathcal E\)-internal presheaf on \(L\) back by
\(f^*\) gives an \(\mathcal F\)-internal presheaf on \(f^*L\).
This base-change functor is left exact and has a right adjoint.
Explicitly, on underlying families its right adjoint is the
slice right adjoint obtained by pulling \(f_*Y\) back along
\(L\to f_*f^*L\); the action is the adjoint transpose of the
action on \(Y\), after the counit. The action equations follow
by transposition, and the same bijection restricts to action
preserving maps. Thus it is also the right adjoint on
presheaves, not merely on underlying objects.
Compose this base-change functor with \(G\).

The composite right adjoint takes values in sheaves.
Indeed, a matching family on a basic cover corresponds, by
the displayed adjunctions, to compatible maps on the
subterminals \(U_v\). Their union is \(U_u\), since \(\chi\)
preserves the finite join defining the cover.
Such maps glue uniquely by descent along their epimorphic
union. This proves exactly the sheaf condition.
Consequently the composite left adjoint, restricted to
sheaves, is left adjoint to this sheaf-valued right adjoint
and remains left exact. It is the inverse image of a
geometric morphism \(g:\mathcal F\to\mathcal S_A\).

The composite sends a constant presheaf with value \(B\)
to \(f^*B\): a value has its representative at \(u=1\),
and restrictions do not change that value.
It also sends a presheaf and its sheafification to the
same object, since its right adjoint is sheaf valued.
Hence \(g^*p^*\cong f^*\), with the compatible adjunction
data, so \(g\) lies over the specified \(f\).
Applied to \(y(u)\), the germ construction yields \(U_u\):
a representative below \(u\) maps into \(U_u\), and an
element of \(U_u\) has its representative at \(u\).

\emph{Recovery and uniqueness.}
Conversely, a geometric morphism over \(f\) sends the
subterminal family \(y(u)\) to a subterminal family indexed
by \(f^*L\). Finite limits and the jointly epimorphic finite
covers show that its classifiers form a bounded lattice
homomorphism \(\chi\).
The preceding construction recovers those subterminals.

It recovers the entire inverse-image functor, not just its
values on opens. Every internal presheaf has its canonical
presentation by representatives \((u,x)\) and the two maps
identifying each representative with its restrictions.
This is an internal dependent-sum and coequalizer
presentation, of the same form as the internal Kan-extension
construction. Sheafifying gives such a presentation of each
sheaf. A geometric inverse image preserves its coequalizer,
finite products, and the dependent sums over objects
transported from the base. Its value is therefore the
germ quotient just constructed.
The resulting natural isomorphism is unique when compatible
with the base and these presentations: its components on
the subterminals are unique, and the presentation maps are
epimorphic. This proves classification up to compatible
isomorphism.

\emph{The generic prime predicate.}
On the site, take the subpresheaf of the constant ring
whose elements at \(u\) are those \(b\in A\) with \(u\leq d(b)\).
Sheafification gives \(S\hookrightarrow p^*A\).
Membership of \(b\) has support \(y(d(b))\).
The lattice equations therefore give the four prime
conditions: top for one, bottom for zero, intersection for
multiplication, and the indicated finite union for addition.
In particular the addition alternative is sheaf-local.
The presentation of \(L_A\), together with
\(f^*L_A\cong L_{f^*A}\), identifies the homomorphisms \(\chi\)
with exactly the prime predicates on \(f^*A\).
The recovered support of each generator says precisely that
this predicate is \(g^*S\).
\end{proof}
\end{annexitem}

\begin{annexitem}{annex:p10-corrected-global-spectrum}{thm:p10-corrected-global-spectrum}{Universality}
\begin{proof}
The preceding theorem constructs \(p\) and the generic prime
subobject \(S\hookrightarrow p^*A\).
By \cref{thm:p10-internal-prime-locality},
\(\mathcal O=S^{-1}(p^*A)\) is a local ring object.

Given \(f\), \(B\), and \(h\) as in the statement, let
\(S_h\hookrightarrow f^*A\) be the inverse image under \(h\)
of the units of \(B\).
The same locality theorem proves that \(S_h\) is a prime
predicate. The site classification supplies \(g\) over \(f\)
with \(g^*S=S_h\), up to the specified canonical isomorphisms.
Preservation of localization by inverse image gives
\[
 g^*\mathcal O\cong S_h^{-1}(f^*A).
\]
The universal property of fractions supplies the unique
extension \(k:g^*\mathcal O\to B\) of \(h\).
The unit criterion proves that \(k\) reflects units, so it
is a local homomorphism.

For uniqueness, consider any factorization of \(h\) through
a local homomorphism from the inverse image of \(\mathcal O\).
The inverse image of the units of \(\mathcal O\) along
\(p^*A\to\mathcal O\) is exactly \(S\), by the numerator
criterion. Since the second homomorphism reflects units,
pulling this statement back identifies the resulting prime
predicate on \(f^*A\) with \(S_h\).
The classification theorem forces the geometric factor to
be the same, up to its unique compatible isomorphism.
Once that factor is fixed, the universal property of the
localization forces its ring homomorphism to be \(k\).
These statements give the asserted unique factorization.
\end{proof}
\end{annexitem}
